% GENERATED FILE. Edit canonical lessons, then run npm run generate:books.
% canonical-source-hash: fnv1a64:761c115e2c14ef7f
% canonical-lessons: MR-C486-ija, MR-C486-khapli, MR-C486-kharcatne, MR-C486-murgalne, MR-C486-sujne

\chapter{Small Injuries}
\label{ch:mr-small-injuries}
\hlchaptermodality{486}

\begin{chapteropening}
\textbf{By the end of this chapter:} \emph{I can talk about an injury, a scab, a graze, a sprain and swelling.}

Complete the last lesson of chapter 486: I can talk about an injury, a scab, a graze, a sprain and swelling.
\end{chapteropening}

\section[\texorpdfstring{\mr{इजा} (ijā)}{इजा (ijā)}]{\mr{इजा} (ijā) — an injury; harm}

\label{lesson:MR-C486-ija}



\begin{quote}
Before the new one: say the Marathi for jaundice, then the Marathi for chickenpox.
\end{quote}

\subsection*{You'll want to know: \mr{इजा}}
\textbf{\mr{इजा}} — \emph{ijā} — \textquotedblleft{}an injury; harm\textquotedblright{}.

\begin{quote}
\mr{अपघातात} \mr{त्याला} \mr{किरकोळ} \mr{इजा} \mr{झाली}.
\end{quote}

\begin{quote}
\emph{apghātāt tyālā kirkoḷ ijā jhālī.}
\end{quote}

\begin{quote}
He was slightly injured in the accident.
\end{quote}

Feminine.

\mr{इजा} \mr{होणे}, to be hurt, with the dative; \mr{इजा} \mr{करणे}, to harm.

\begin{cousinweb}
\textbf{\mr{इजा}} (\emph{ijā}) is a loan from Persian \emph{īzā}, \textquotedblleft{}harm, injury\textquotedblright{}, usually traced to Arabic; Marathi writes the \emph{z} as \mr{ज}.
\end{cousinweb}

\begin{grammarlens}[title={What you've built}]
One more word for describing what you lived and felt.
\end{grammarlens}

\subsection*{Your turn}
\begin{itemize}
\raggedright
  \item \emph{Say it:} \emph{ijā}
  \item \emph{Say it:} \emph{ijā}, once more
  \item \emph{Say it:} \emph{apghātāt tyālā kirkoḷ ijā jhālī.} — He was slightly injured in the accident.
  \item \emph{Recall:} say the Marathi for jaundice, then the Marathi for chickenpox, then say \emph{ijā} again
\end{itemize}

\begin{tcolorbox}[breakable,colback=teal!4,colframe=teal!35!black,title={Before you move on}]
What does \mr{इजा} mean? (\textquotedblleft{}An injury; harm\textquotedblright{}.) Say it once more.
\end{tcolorbox}
\section[\texorpdfstring{\mr{खपली} (khaplī)}{खपली (khaplī)}]{\mr{खपली} (khaplī) — a scab}

\label{lesson:MR-C486-khapli}



\begin{quote}
Before the new one: say the Marathi for chickenpox, then the Marathi for an injury.
\end{quote}

\subsection*{You'll want to know: \mr{खपली}}
\textbf{\mr{खपली}} — \emph{khaplī} — \textquotedblleft{}a scab\textquotedblright{}.

\begin{quote}
\mr{जखमेवर} \mr{आता} \mr{खपली} \mr{आली} \mr{आहे}.
\end{quote}

\begin{quote}
\emph{jakhmevar ātā khaplī ālī āhe.}
\end{quote}

\begin{quote}
A scab has formed on the wound now.
\end{quote}

Feminine.

\mr{खपली} \mr{धरणे}, a scab forms; \mr{खपली} \mr{काढणे}, to pick a scab.

\begin{cousinweb}
\textbf{\mr{खपली}} (\emph{khaplī}) is a Marathi noun whose origin is not settled.
\end{cousinweb}

\begin{grammarlens}[title={What you've built}]
One more word for describing what you lived and felt.
\end{grammarlens}

\subsection*{Your turn}
\begin{itemize}
\raggedright
  \item \emph{Say it:} \emph{khaplī}
  \item \emph{Say it:} \emph{khaplī}, once more
  \item \emph{Say it:} \emph{jakhmevar ātā khaplī ālī āhe.} — A scab has formed on the wound now.
  \item \emph{Recall:} say the Marathi for chickenpox, then the Marathi for an injury, then say \emph{khaplī} again
\end{itemize}

\begin{tcolorbox}[breakable,colback=teal!4,colframe=teal!35!black,title={Before you move on}]
What does \mr{खपली} mean? (\textquotedblleft{}A scab\textquotedblright{}.) Say it once more.
\end{tcolorbox}
\section[\texorpdfstring{\mr{खरचटणे} (kharcaṭṇe)}{खरचटणे (kharcaṭṇe)}]{\mr{खरचटणे} (kharcaṭṇe) — to get grazed}

\label{lesson:MR-C486-kharcatne}



\begin{quote}
Before the new one: say the Marathi for an injury, then the Marathi for a scab.
\end{quote}

\subsection*{You'll want to know: \mr{खरचटणे}}
\textbf{\mr{खरचटणे}} — \emph{kharcaṭṇe} — \textquotedblleft{}to get grazed\textquotedblright{}.

\begin{quote}
\mr{सायकलवरून} \mr{पडल्यावर} \mr{त्याच्या} \mr{गुडघ्याला} \mr{खरचटलं}.
\end{quote}

\begin{quote}
\emph{sāykalvarūn paḍlyāvar tyācyā guḍghyālā kharcaṭlaṁ.}
\end{quote}

\begin{quote}
When he fell off his bicycle he grazed his knee.
\end{quote}

Past (to him …): \mr{त्याला} \mr{खरचटलं}, \emph{tyālā kharcaṭlaṁ}.

Impersonal: the grazed part takes -\mr{ला}: \mr{हाताला} \mr{खरचटलं}.

\begin{cousinweb}
\textbf{\mr{खरचटणे}} (\emph{kharcaṭṇe}) is a Marathi verb whose origin is not settled.
\end{cousinweb}

\begin{grammarlens}[title={What you've built}]
One more word for describing what you lived and felt.
\end{grammarlens}

\subsection*{Your turn}
\begin{itemize}
\raggedright
  \item \emph{Say it:} \emph{kharcaṭṇe}
  \item \emph{Say it:} \emph{kharcaṭṇe}, once more
  \item \emph{Say it:} \emph{sāykalvarūn paḍlyāvar tyācyā guḍghyālā kharcaṭlaṁ.} — When he fell off his bicycle he grazed his knee.
  \item \emph{Recall:} say the Marathi for an injury, then the Marathi for a scab, then say \emph{kharcaṭṇe} again
\end{itemize}

\begin{tcolorbox}[breakable,colback=teal!4,colframe=teal!35!black,title={Before you move on}]
What does \mr{खरचटणे} mean? (\textquotedblleft{}To get grazed\textquotedblright{}.) Say it once more.
\end{tcolorbox}
\section[\texorpdfstring{\mr{मुरगळणे} (murgaḷṇe)}{मुरगळणे (murgaḷṇe)}]{\mr{मुरगळणे} (murgaḷṇe) — to sprain; to twist}

\label{lesson:MR-C486-murgalne}



\begin{quote}
Before the new one: say the Marathi for a scab, then the Marathi for to get grazed.
\end{quote}

\subsection*{You'll want to know: \mr{मुरगळणे}}
\textbf{\mr{मुरगळणे}} — \emph{murgaḷṇe} — \textquotedblleft{}to sprain; to twist\textquotedblright{}.

\begin{quote}
\mr{खेळताना} \mr{माझा} \mr{घोटा} \mr{मुरगळला}.
\end{quote}

\begin{quote}
\emph{kheḷtānā mājhā ghoṭā murgaḷlā.}
\end{quote}

\begin{quote}
I sprained my ankle while playing.
\end{quote}

Past, with -\mr{ने} (he …): \mr{त्याने} \mr{कान} \mr{मुरगळला}, \emph{tyāne kān murgaḷlā}.

Of a sprain the joint is the subject: \mr{माझा} \mr{पाय} \mr{मुरगळला}. With -\mr{ने} it is to twist: \mr{कान} \mr{मुरगळणे}.

\begin{cousinweb}
\textbf{\mr{मुरगळणे}} (\emph{murgaḷṇe}) is a Marathi verb whose origin is not settled.
\end{cousinweb}

\begin{grammarlens}[title={What you've built}]
One more word for describing what you lived and felt.
\end{grammarlens}

\subsection*{Your turn}
\begin{itemize}
\raggedright
  \item \emph{Say it:} \emph{murgaḷṇe}
  \item \emph{Say it:} \emph{murgaḷṇe}, once more
  \item \emph{Say it:} \emph{kheḷtānā mājhā ghoṭā murgaḷlā.} — I sprained my ankle while playing.
  \item \emph{Recall:} say the Marathi for a scab, then the Marathi for to get grazed, then say \emph{murgaḷṇe} again
\end{itemize}

\begin{tcolorbox}[breakable,colback=teal!4,colframe=teal!35!black,title={Before you move on}]
What does \mr{मुरगळणे} mean? (\textquotedblleft{}To sprain; to twist\textquotedblright{}.) Say it once more.
\end{tcolorbox}
\section[\texorpdfstring{\mr{सुजणे} (sujṇe)}{सुजणे (sujṇe)}]{\mr{सुजणे} (sujṇe) — to swell}

\label{lesson:MR-C486-sujne}



\begin{quote}
Before the new one: say the Marathi for to get grazed, then the Marathi for to sprain.
\end{quote}

\subsection*{You'll want to know: \mr{सुजणे}}
\textbf{\mr{सुजणे}} — \emph{sujṇe} — \textquotedblleft{}to swell\textquotedblright{}.

\begin{quote}
\mr{मधमाशी} \mr{चावल्यामुळे} \mr{त्याचा} \mr{हात} \mr{सुजला}.
\end{quote}

\begin{quote}
\emph{madhmāśī cāvlyāmuḷe tyācā hāt sujlā.}
\end{quote}

\begin{quote}
His hand swelled up from a bee sting.
\end{quote}

Past (he …): \mr{तो} \mr{सुजला}, \emph{to sujlā}.

Usually the swollen part is the subject: \mr{त्याचा} \mr{पाय} \mr{सुजला}. \mr{सूज}, swelling: \mr{सूज} \mr{येणे}.

\begin{cousinweb}
\textbf{\mr{सुजणे}} (\emph{sujṇe}) is a tadbhava, inherited from Sanskrit \emph{śūyate}, \textquotedblleft{}is swollen\textquotedblright{}; CDIAL lists Marathi \emph{sujṇẽ}, \textquotedblleft{}to swell\textquotedblright{} (CDIAL 12568).
\end{cousinweb}

\begin{grammarlens}[title={What you've built}]
One more word for describing what you lived and felt.
\end{grammarlens}

\subsection*{Your turn}
\begin{itemize}
\raggedright
  \item \emph{Say it:} \emph{sujṇe}
  \item \emph{Say it:} \emph{sujṇe}, once more
  \item \emph{Say it:} \emph{madhmāśī cāvlyāmuḷe tyācā hāt sujlā.} — His hand swelled up from a bee sting.
  \item \emph{Recall:} say the Marathi for to get grazed, then the Marathi for to sprain, then say \emph{sujṇe} again
\end{itemize}

\begin{tcolorbox}[breakable,colback=teal!4,colframe=teal!35!black,title={Before you move on}]
What does \mr{सुजणे} mean? (\textquotedblleft{}To swell\textquotedblright{}.) Say it once more.
\end{tcolorbox}
